% !TeX root = ../main.tex
% ---------------------------------------------------------------------------
% Shafts, Portals and Access Structures
% Scope extension: not derived from the RSE3010 weekly material.
% Key sources: BTS/ICE Specification for Tunnelling; Puller, Deep Excavations;
% Peck (1969) on braced excavations; ITA WG on shafts.
% ---------------------------------------------------------------------------
\chapter{Shafts, Portals and Access Structures}
\label{ch:shafts-portals}


\section{Role and construction methods}
\label{sec:role-and-construction-methods}

\subsection{The role of shafts, portals and adits}
\label{sec:access-role}
\scopenote{Access, ventilation, drainage, launch and reception, and permanent
operational function. Why these structures carry disproportionate risk and
cost per metre.}

\subsection{Shaft types and sinking methods}
\label{sec:shaft-methods}
\scopenote{Drill and blast sinking, mechanical excavation, caisson sinking,
slurry wall shafts, secant and contiguous pile shafts, and shaft boring
machines. Selection by depth, diameter, ground and water.}

\section{Shaft stress and lining design}
\label{sec:shaft-stress-and-lining-design}

\subsection{Stress around a vertical opening}
\label{sec:shaft-stress}
\scopenote{Stress concentration around a circular shaft in a horizontally
anisotropic field, the difference from a horizontal tunnel, and the
implications for lining loads and spalling.}

\subsection{Circular shaft lining design}
\label{sec:shaft-lining}
\scopenote{Hoop compression as the efficient load path, ring design for ground
and water pressure, and the sensitivity of a compression ring to construction
out-of-roundness.}

\section{Sinking methods}
\label{sec:sinking-methods}

\subsection{Shaft sinking in soft ground}
\label{sec:soft-ground-shafts}
\scopenote{Diaphragm walls, secant piles and caissons; base heave and hydraulic
uplift; base plugs and jet grouted plugs; dewatering versus cut-off.}

\subsection{Raise boring and blind boring}
\label{sec:raise-boring}
\scopenote{Raise boring for ventilation and penstock shafts, pilot hole
accuracy, reaming, and the stability requirements for an unsupported bore.}

\section{Portals}
\label{sec:portals-group}

\subsection{Portal design and slope stability}
\label{sec:portals}
\scopenote{The portal as a combined slope and tunnel problem; face and side
slope stability, kinematic assessment, and the low-cover first rounds where
most portal failures occur.}

\subsection{Portal support and canopies}
\label{sec:portal-support}
\scopenote{Pipe canopies, forepoling, soil nailing, shotcrete facing and
retaining structures; sequencing the transition from open cut to bored
tunnel.}

\section{Openings and deep excavations}
\label{sec:openings-and-deep-excavations}

\subsection{Adits, cross-passages and niches}
\label{sec:cross-passages}
\scopenote{Openings in an existing lining: three-dimensional stress
redistribution, opening reinforcement, and the ground treatment usually
required before breakout.}

\subsection{Deep excavations and station boxes}
\label{sec:station-boxes}
\scopenote{Retaining wall selection, propping and anchoring, top-down
construction, base stability, and wall deflection versus ground movement
relationships. Damage assessment in Chapter~\ref{ch:settlement-damage}.}

\subsection{Interfaces and transitions}
\label{sec:interfaces}
\scopenote{Bored-to-cut-and-cover transitions, differential stiffness and
settlement at interfaces, and waterproofing continuity across them.}
